\begin{solution}{85}
    \begin{subsolution}
        \begin{subsubsolution}
            设大气压强为 $P_{0}$，毛细管的内径为 $r$，外径为 $R$，水的密度为 $\rho_{0}$，表面张力系数为 $\sigma$，重力加速度为 $g$。对于水柱的上端，水与管壁完全浸润，接触角为零 $\theta=0$，水柱上端内表面压强 $P_{1}$ 为
            \begin{equation}
                P_{1} = P_{0} - \frac{2\sigma}{r}
                \label{eq:1.s85}
            \end{equation}
            对于水柱下端，同理可知，表面压强 $P_{2}$ 为
            \begin{equation}
                P_{2} = P_{0} + \frac{2\sigma}{a}
                \label{eq:2.s85}
            \end{equation}
            水柱平衡时有
            \begin{equation}
                P_{2} = P_{1} + \rho_{0} g h
                \label{eq:3.s85}
            \end{equation}
            由\eqref{eq:1.s85}--\eqref{eq:3.s85}式得
            \begin{equation}
                a = \frac{2\sigma r}{\rho_{0} g h r - 2\sigma}
                \label{eq:4.s85}
            \end{equation}
            由\eqref{eq:4.s85}式和题给数据得
            \begin{equation}
                a = 2.79 \times 10^{-3}\,\mathrm{m}
                \label{eq:5.s85}
            \end{equation}
        \end{subsubsolution}
        \begin{subsubsolution}
            将该毛细管长度的三分之一竖直浸入水中后，设管内水柱顶部与管外水平的水面高度差为 $h_{1}$，由平衡条件得
            \begin{equation}
                \rho_{0} \cdot \pi r^{2} h_{1} g = 2\pi r \cdot \sigma
                \label{eq:6.s85}
            \end{equation}
            即
            \[
                h_{1} = \frac{2\sigma}{\rho_{0} g r}
            \]
            由上式和题给数据得
            \begin{equation}
                h_{1} = 2.97\,\mathrm{cm} < 4.00\,\mathrm{cm}
                \label{eq:7.s85}
            \end{equation}
            毛细管内壁所受的附着力 $F_{1}$ 大小等于管外水平的水面以上管内水柱的重力大小 $G_{1}$，即
            \begin{equation}
                F_{1} = G_{1} = \rho_{0} g \pi r^{2} h_{1}
                \label{eq:8.s85}
            \end{equation}
            方向向下。玻璃管外壁所受的附着力 $F_{2}$ 大小为
            \begin{equation}
                F_{2} = \sigma 2\pi R
                \label{eq:9.s85}
            \end{equation}
            方向向下。设玻璃管长度为 $h$，玻璃管所受的重力大小为
            \begin{equation}
                G = 2\rho_{0} g \pi (R^{2} - r^{2}) h
                \label{eq:10.s85}
            \end{equation}
            方向向下。玻璃管所受到的浮力大小 $F_{0}$ 为
            \begin{equation}
                F_{0} = \frac{1}{3} h \rho_{0} g \pi (R^{2} - r^{2})
                \label{eq:11.s85}
            \end{equation}
            方向向上。设向上的拉力大小为 $F$，由力平衡条件得
            \begin{equation}
                F = G + G_{1} + F_{2} - F_{0}
                \label{eq:12.s85}
            \end{equation}
            由\eqref{eq:7.s85}--\eqref{eq:12.s85}式得
            \begin{align}
                F &= \rho g \pi (R^{2} - r^{2}) h + \rho_{0} g \pi r^{2} h_{1} + \sigma 2\pi R - \frac{1}{3} \rho_{0} g \pi (R^{2} - r^{2}) h \notag\\
                &= \frac{5}{3} \rho_{0} g \pi h (R^{2} - r^{2}) + 2\pi \sigma (R+r)
                \label{eq:13.s85}
            \end{align}
            由\eqref{eq:13.s85}式和题给数据得
            \begin{equation}
                F = 7.87 \times 10^{-2}\,\mathrm{N}
                \label{eq:14.s85}
            \end{equation}
        \end{subsubsolution}
    \end{subsolution}
\end{solution}